% Préambule commun des schémas TikZ — charte Liouaï
% Chaque schéma commence par % Préambule commun des schémas TikZ — charte Liouaï
% Chaque schéma commence par % Préambule commun des schémas TikZ — charte Liouaï
% Chaque schéma commence par \input{_preambule} puis \begin{document}.
\documentclass[tikz,border=6pt]{standalone}
\usepackage[T1]{fontenc}
\usepackage[utf8]{inputenc}
\usepackage[scaled=0.95]{helvet}
\renewcommand{\familydefault}{\sfdefault}
\usepackage{amsmath,amssymb}
\usepackage[european]{circuitikz}
\usetikzlibrary{arrows.meta,positioning,calc,angles,quotes,decorations.pathreplacing,backgrounds}

% Couleurs du logo Liouaï Digital Education
\definecolor{navy}{HTML}{243743}
\definecolor{vert}{HTML}{28B78D}
\definecolor{vertfonce}{HTML}{167A5C}
\definecolor{vertclair}{HTML}{EAF7F2}
\definecolor{gris}{HTML}{A1A9AE}
\definecolor{gristexte}{HTML}{5C6870}
\definecolor{fond}{HTML}{F3F5F6}
\definecolor{jaune}{HTML}{E5C800}
\definecolor{jauneclair}{HTML}{FFFBD6}

\tikzset{
  every picture/.style={line cap=round, line join=round, color=navy, text=navy},
  fleche/.style={-{Stealth[length=2.6mm,width=2mm]}, line width=1.1pt},
  vecteur/.style={fleche, line width=1.4pt},
  titre/.style={font=\bfseries\small, text=navy},
  note/.style={font=\footnotesize, text=gristexte},
  boite/.style={draw=navy, line width=0.9pt, rounded corners=5pt, fill=fond, align=center, inner sep=6pt},
}
 puis \begin{document}.
\documentclass[tikz,border=6pt]{standalone}
\usepackage[T1]{fontenc}
\usepackage[utf8]{inputenc}
\usepackage[scaled=0.95]{helvet}
\renewcommand{\familydefault}{\sfdefault}
\usepackage{amsmath,amssymb}
\usepackage[european]{circuitikz}
\usetikzlibrary{arrows.meta,positioning,calc,angles,quotes,decorations.pathreplacing,backgrounds}

% Couleurs du logo Liouaï Digital Education
\definecolor{navy}{HTML}{243743}
\definecolor{vert}{HTML}{28B78D}
\definecolor{vertfonce}{HTML}{167A5C}
\definecolor{vertclair}{HTML}{EAF7F2}
\definecolor{gris}{HTML}{A1A9AE}
\definecolor{gristexte}{HTML}{5C6870}
\definecolor{fond}{HTML}{F3F5F6}
\definecolor{jaune}{HTML}{E5C800}
\definecolor{jauneclair}{HTML}{FFFBD6}

\tikzset{
  every picture/.style={line cap=round, line join=round, color=navy, text=navy},
  fleche/.style={-{Stealth[length=2.6mm,width=2mm]}, line width=1.1pt},
  vecteur/.style={fleche, line width=1.4pt},
  titre/.style={font=\bfseries\small, text=navy},
  note/.style={font=\footnotesize, text=gristexte},
  boite/.style={draw=navy, line width=0.9pt, rounded corners=5pt, fill=fond, align=center, inner sep=6pt},
}
 puis \begin{document}.
\documentclass[tikz,border=6pt]{standalone}
\usepackage[T1]{fontenc}
\usepackage[utf8]{inputenc}
\usepackage[scaled=0.95]{helvet}
\renewcommand{\familydefault}{\sfdefault}
\usepackage{amsmath,amssymb}
\usepackage[european]{circuitikz}
\usetikzlibrary{arrows.meta,positioning,calc,angles,quotes,decorations.pathreplacing,backgrounds}

% Couleurs du logo Liouaï Digital Education
\definecolor{navy}{HTML}{243743}
\definecolor{vert}{HTML}{28B78D}
\definecolor{vertfonce}{HTML}{167A5C}
\definecolor{vertclair}{HTML}{EAF7F2}
\definecolor{gris}{HTML}{A1A9AE}
\definecolor{gristexte}{HTML}{5C6870}
\definecolor{fond}{HTML}{F3F5F6}
\definecolor{jaune}{HTML}{E5C800}
\definecolor{jauneclair}{HTML}{FFFBD6}

\tikzset{
  every picture/.style={line cap=round, line join=round, color=navy, text=navy},
  fleche/.style={-{Stealth[length=2.6mm,width=2mm]}, line width=1.1pt},
  vecteur/.style={fleche, line width=1.4pt},
  titre/.style={font=\bfseries\small, text=navy},
  note/.style={font=\footnotesize, text=gristexte},
  boite/.style={draw=navy, line width=0.9pt, rounded corners=5pt, fill=fond, align=center, inner sep=6pt},
}

\begin{document}
\begin{tikzpicture}[x=0.75cm, y=0.6cm]
\newcommand{\grille}[2]{%
  \foreach \i in {0,...,3} \foreach \j in {0,...,2}
    \draw[gris, fill=fond] ({#1+\i}, {#2-\j}) rectangle ++(1,-1);}
% --- axis=0 : on écrase les lignes, un total par colonne
\node[titre] at (2, 2.2) {\texttt{D.sum(axis=0)}};
\node[note] at (2, 1.3) {additionne en descendant};
\grille{0}{0}
\foreach \i in {0,...,3} \draw[fleche, vertfonce] (\i+0.5, -0.2) -- (\i+0.5, -2.8);
\foreach \i in {0,...,3} \draw[draw=vertfonce, fill=vert] (\i, -3.6) rectangle ++(1,-1);
\node[note, align=center] at (2, -5.4) {un total par colonne\\ (4 nombres : par catégorie)};
\node[note, rotate=90] at (-0.6, -1.5) {mois};
\node[note] at (2, 0.45) {4 catégories};
% --- axis=1 : on écrase les colonnes, un total par ligne
\begin{scope}[xshift=7.2cm]
\node[titre] at (2.5, 2.2) {\texttt{D.sum(axis=1)}};
\node[note] at (2.5, 1.3) {additionne le long de chaque ligne};
\grille{0}{0}
\foreach \j in {0,...,2} \draw[fleche, vertfonce] (0.2, -\j-0.5) -- (3.8, -\j-0.5);
\foreach \j in {0,...,2} \draw[draw=vertfonce, fill=vert] (4.6, -\j) rectangle ++(1,-1);
\node[note, align=center] at (2.8, -5.4) {un total par ligne\\ (12 nombres : par mois)};
\node[note] at (2, 0.45) {4 catégories};
\end{scope}
\end{tikzpicture}
\end{document}
